% C-24 — Conclusions' basal-ganglia mechanism paragraph, redundancy pass
% Origin: intake/pending/redundancy_closing_audit.md (R2-01-derivative redundancy audit).
% Previously flagged as an unselected option during C-12 ("condense conclusions.tex's
% STN/GPe/GPi paragraph vs methodology.tex") -- executed now as part of this audit.
% File: sections/conclusions.tex

% --- BEFORE ---
% In scenarios presenting simultaneous appetitive and aversive stimuli, the robot consistently
% prioritized the most relevant objective through the basal-ganglia-inspired arbitration module.
% The interaction among the STR, STN, GPe, and GPi implemented a winner-take-all mechanism,
% actively suppressing responses toward non-priority stimuli. During the tests, predominant
% activation in the channel associated with the selected stimulus effectively inhibited
% competing channels, reducing the probability of erratic trajectory changes. This behavior
% yielded direction-switching latencies below 600~ms, ensuring fast and stable responses under
% sensory conflict. The alignment between decision-neuron activity peaks and executed maneuvers
% supports the hypothesis that a basal-ganglia-inspired architecture combined with multimodal
% sensory processing offers clear advantages in environments where goal prioritization is
% critical for mission success.

% --- AFTER ---
In scenarios presenting simultaneous appetitive and aversive stimuli, the basal-ganglia-inspired
winner-take-all arbitration module (Section~\ref{sssec:Basal_Ganglia_Module}) \revblue{consistently
prioritized the most relevant objective, its predominant channel activation inhibiting competing
channels and reducing the probability of erratic trajectory changes; this behavior yielded
direction-switching latencies below 600~ms.} The alignment between decision-neuron activity peaks
and executed maneuvers supports the hypothesis that a basal-ganglia-inspired architecture combined
with multimodal sensory processing offers clear advantages in environments where goal
prioritization is critical for mission success.

% Also added \label{sssec:Basal_Ganglia_Module} to methodology.tex's
% \subsubsection{Basal Ganglia Module} (previously unlabeled) so the cross-reference above
% resolves.

% Rationale: the "STR/STN/GPe/GPi implemented a winner-take-all mechanism, actively suppressing
% responses toward non-priority stimuli" sentence re-explains the same architecture already
% described in methodology.tex's Basal Ganglia Module subsubsection ("The GPe inhibits
% lower-intensity behaviors... the GPi complements this mechanism by suppressing competing
% responses to ensure that only a single action is executed"). Replaced the mechanism
% re-explanation with a cross-reference, kept the results-specific content (behavioral outcome,
% 600ms latency figure, closing hypothesis sentence) untouched.
% NOTE: the "600ms" figure itself is not independently re-verified by this patch -- it was
% already flagged elsewhere (intake/pending/redundancy_closing_audit.md's C-12 discussion) as an
% open traceability question, separate from this redundancy fix.
